\subsection{Jenis Testing yang Dilakukan}

ResQfood menjalani berbagai jenis pengujian untuk memastikan kualitas dan keandalan aplikasi:

\begin{table}[H]
\centering
\caption{Jenis Pengujian ResQfood}
\label{tab:testing}
\begin{tabularx}{\textwidth}{lXX}
\toprule
\textbf{Jenis Testing} & \textbf{Deskripsi} & \textbf{Tools} \\
\midrule
Unit Testing & Pengujian fungsi individual & Jest, Vitest \\
Integration Testing & Pengujian integrasi antar modul & Jest, Supertest \\
API Testing & Pengujian endpoint API & Postman, Supertest \\
UI Testing & Pengujian antarmuka pengguna & Cypress \\
Usability Testing & Pengujian kemudahan penggunaan & Manual testing \\
Performance Testing & Pengujian performa dan beban & Artillery \\
Security Testing & Pengujian keamanan dasar & OWASP ZAP \\
\bottomrule
\end{tabularx}
\end{table}

\subsection{Unit Testing}

\subsubsection{Contoh Test Case: Food Eligibility Score}
\begin{verbatim}
describe('Food Eligibility Score', () => {
  test('should return score > 80 for fresh food', () => {
    const foodData = {
      productionTime: new Date(),
      shelfLife: 24,
      packagingCondition: 'good',
      storageTemp: 4
    };
    const score = calculateFES(foodData);
    expect(score).toBeGreaterThan(80);
  });

  test('should return score < 50 for expired food', () => {
    const foodData = {
      productionTime: new Date('2024-01-01'),
      shelfLife: 2,
      packagingCondition: 'poor',
      storageTemp: 15
    };
    const score = calculateFES(foodData);
    expect(score).toBeLessThan(50);
  });
});
\end{verbatim}

\subsection{API Testing}

\begin{table}[H]
\centering
\caption{Hasil Pengujian API}
\label{tab:api_testing}
\begin{tabularx}{\textwidth}{lXX}
\toprule
\textbf{Endpoint} & \textbf{Status} & \textbf{Response Time} \\
\midrule
POST /api/auth/register & 201 Created & $<$ 200ms \\
POST /api/auth/login & 200 OK & $<$ 150ms \\
GET /api/products & 200 OK & $<$ 100ms \\
POST /api/products & 201 Created & $<$ 250ms \\
POST /api/orders & 201 Created & $<$ 300ms \\
PUT /api/orders/:id/status & 200 OK & $<$ 150ms \\
\bottomrule
\end{tabularx}
\end{table}

\subsection{Usability Testing}

Pengujian kegunaan dilakukan dengan 5 partisipan yang mewakili target pengguna:

\begin{table}[H]
\centering
\caption{Hasil Usability Testing}
\label{tab:usability}
\begin{tabularx}{\textwidth}{lXX}
\toprule
\textbf{Skenario} & \textbf{Skor Kemudahan} & \textbf{Waktu Rata-rata} \\
\midrule
Registrasi akun baru & 4.2/5 & 2 menit \\
Login dan navigasi & 4.5/5 & 1 menit \\
Membeli produk (Customer) & 4.0/5 & 3 menit \\
Menambah produk (UMKM) & 3.8/5 & 4 menit \\
Mengisi Food Eligibility & 3.5/5 & 5 menit \\
\bottomrule
\end{tabularx}
\end{table}

\subsection{Performance Testing}

\begin{table}[H]
\centering
\caption{Hasil Pengujian Performa}
\label{tab:performance}
\begin{tabularx}{\textwidth}{lXX}
\toprule
\textbf{Metrik} & \textbf{Target} & \textbf{Hasil} \\
\midrule
Response Time (rata-rata) & $<$ 500ms & 180ms \\
Response Time (p95) & $<$ 1000ms & 450ms \\
Concurrent Users & 100 & 150 \\
Error Rate & $<$ 1\% & 0.2\% \\
Uptime & 99\% & 99.5\% \\
\bottomrule
\end{tabularx}
\end{table}

\subsection{Limitasi Pengujian}

\begin{itemize}[leftmargin=*]
    \item Pengujian keamanan baru dilakukan pada level dasar (OWASP Top 10). Pengujian penetrasi lengkap membutuhkan waktu dan ahli keamanan khusus.
    \item Pengujian performa dilakukan pada skala kecil (150 pengguna simultan). Pengujian skala besar membutuhkan infrastruktur production-ready.
    \item Fitur prediksi produksi belum dapat diuji secara menyeluruh karena membutuhkan data historis yang cukup.
\end{itemize}
